\section{Troubleshooting}
\label{sec:troubleshooting}

\subsection{``Cannot start analysis''}

A modal dialog with one of these messages:

\begin{description}[leftmargin=1.2cm,style=nextline]
  \item[Need at least 2 images]
    Drop an image folder into the Images panel.
  \item[Define a Region of Interest first]
    Use the Region of Interest toolbar to draw or import one on
    frame 1.
  \item[Region of Interest is empty]
    You drew something, then clicked Cut over the whole region.
    Click Clear or redraw.
  \item[Region of Interest too small]
    The bounding box is smaller than $2 \times \text{Subset Step}$.
    Enlarge the Region or reduce Subset Step.
\end{description}

\subsection{``DIC analysis failed''}

Shows an exception type and message. The full traceback stays in
the console log --- scroll up to find it. Common causes:
\begin{itemize}
  \item Images of different sizes (should be caught at load; if you
        see it at Run, report a bug).
  \item Extremely low subset correlation (black-on-black image or
        no speckle). Check your input data first.
\end{itemize}

\subsection{Slow run}

\begin{itemize}
  \item Subset Step too small: halve the step $\to$ 4$\times$ node
        count $\to$ 4$\times$ runtime. Start with step = 16 for
        large images.
  \item AL-DIC solver: $\sim 3\times$ slower than Local DIC.
        Preview with Local DIC, switch to AL-DIC for the final run.
  \item Search Range much larger than actual motion: every FFT call
        costs proportional to (search)$^2$. Set it just above your
        expected inter-frame motion.
\end{itemize}

\subsection{Poor results in specific regions}

\begin{itemize}
  \item Noisy or fuzzy edges of the Region of Interest: enable
        \emph{Refine Outer Boundary} to densify mesh near the edge.
  \item Hole boundaries (e.g.\ bubble edges): enable \emph{Refine
        Inner Boundary}.
  \item Whole regions that simply failed to converge: check the
        console log for ``n\_bad\_pts'' warnings. Usually means
        Search Range was too small or the speckle was washed out
        there.
\end{itemize}

\subsection{Large rotation gives crazy strain values}

You are probably using \emph{Infinitesimal} strain with a rotation
$>$ 2\textdegree. Infinitesimal strain interprets rigid rotation as
strain of magnitude $\cos\theta - 1$ (about $-0.13$ at 30\textdegree).
Switch the Strain window's \emph{Strain type} to
\textbf{Green--Lagrangian}, which is exactly zero under any rigid
rotation.

\subsection{Refine brush grayed out}

Brush is gated to frame 1 only. Click into the image list or use
the frame slider to return to frame 1. The button re-enables
automatically.

\subsection{Session file won't open}

The fatal-error dialog shows the specific reason:
\begin{itemize}
  \item \texttt{not valid JSON}: file corrupted or hand-edited
        wrong. Open it in a text editor, validate against a JSON
        linter.
  \item \texttt{Unsupported schema version}: the file was saved by a
        newer pyALDIC version than the one you are running. Upgrade
        pyALDIC.
  \item \texttt{Image folder ... no longer exists}: this is a
        warning, not a failure. Parameters and Regions of Interest
        still load; re-select the image folder manually.
\end{itemize}
